%% kapa: exam=1 round=30 year=2019 period=1 form=A subject=economics subject_ko=경제학원론 seq=9 number=49
%% source: Q-Net 감정평가사 기출문제 (cst00309) artlSeq=5205205 file=extracted/1008/5205205/2019년 제30회 감정평가사 1차시험 문제지/2019년 제30회 감정평가사 1차시험 문제지/2019년도 제30회 감정평가사 1차 1교시 A형.hwp
%% conversion: kordoc 4.24.51 -> tools/gen1.py (auto); review=proofread-sonnet; flags=none
\begin{question}{49}
보상수요(compensated demand)에 관한 설명으로 옳지 않은 것은?
\begin{choices}
  \item 가격변화에서 대체효과만 고려한 수요개념이다.
  \item 기펜재의 보상수요곡선은 우하향하지 않는다.
  \item 소비자잉여를 측정하는 데 적절한 수요개념이다.
  \item 수직선형태 보상수요곡선의 대체효과는 항상 0이다.
  \item 소득효과가 0이면 통상적 수요(ordinary demand)와 일치한다.
\end{choices}
\end{question}
\answer{2}
